%
% Residuen und die Form des Integrals
%
\subsection{Residuen und die Form des Integrals
\label{buch:ratint:rothsteintrager:subsection:residuen}}
In den komplexen Zahlen lässt sich der quadratfreie Nenner $q(z)$ von
\eqref{buch:ratint:rothsteintrager:eqn:integral}
in ein Produkt verschiedener Linearfaktoren
\begin{equation}
q(z)
=
(z-\alpha_1)\cdots(z-\alpha_n)
\label{buch:ratint:rothsteintrager:eqn:faktorisierung}
\end{equation}
zerlegen.
Da die Nullstellen alle verschieden sind, hat die Partialbruchzerlegung
des Integranden die Form
\begin{equation}
\frac{p(z)}{q(z)}
=
\frac{A_1}{z-\alpha_1}
+
\ldots
+
\frac{A_n}{z-\alpha_n}.
\label{buch:ratint:rothsteintrager:eqn:integrand}
\end{equation}
In dieser Form lässt sich das Integral jedoch sofort angeben, es ist
\begin{equation}
\int
\frac{p(z)}{q(z)}
\,dz
=
A_1\log(z-\alpha_1)
+
\ldots
+
A_n\log(z-\alpha_n).
\label{buch:ratint:rothsteintrager:eqn:integralform}
\end{equation}
Der Ausdruck kann noch etwas vereinfacht werden, wenn zwei der
Koeffizienten gleich sind.
Ist nämlich $A_k=A_l$, dann können die beiden Summanden in
$A_k\log\bigl((z-\alpha_k)(z-\alpha_k)\bigr)$
zusammengefasst werden.
Für entgegengesetzte Koeffizienten $A_k=-A_l$ folgt entsprechend
\[
A_k\log \frac{z-\alpha_k}{z-\alpha_l}.
\]
Diese Form tritt zum Beispiel bei der Integration reeller Funktionen
auf, wo die Nullstellen des Nenners paarweise konjugiert komplex sind.

Wenn die Faktorisierung
\eqref{buch:ratint:rothsteintrager:eqn:faktorisierung}
von $q(z)$ bekannt ist, ist das Integrationsproblem an dieser
Stelle gelöst.
Wie früher bereits dargelegt wurde, kann man jedoch nicht davon
ausgehen, dass das Faktorisierungsproblem gelöst werden kann.
Es müssen also alternative Methoden gesucht werden, die ohne die 
die Faktorisierung auskommen.

%
% Die Koeffizienten A_k und Residuen
%
\subsubsection{Die Koeffizienten $A_k$ und Residuen}
Die Form 
\eqref{buch:ratint:rothsteintrager:eqn:integrand}
des Integranden des logarithmischen Teils zeigt, dass er isolierte
Singularitäten hat und daher um jede Singularität in eine Laurent-Reihe
entwickelt werden kann.
\index{Singularitat@Singularität}%
Man kann sogar sagen, dass die Singularitäten Pole erster Ordnung
sein müssen.

Sei jetzt $\alpha$ eine dieser Pole, dann kann der Integrand
in eine Laurent-Reihe
\index{Laurent-Reihe}%
\[
\frac{p(z)}{q(z)}
=
\frac{b_1}{z-\alpha}
+
a_0+a_1(z-\alpha) + a_2(z-\alpha)^2+\ldots
\]
um den Punkt $\alpha$ entwickelt werden.
Gliedweise Integration der Reihe ergibt in einer Umgebung von $\alpha$
\[
\int
\frac{p(z)}{q(z)}
\,dz
=
b_1\log(z-\alpha)
+
a_0(z-\alpha) + \frac{a_1}{2}(z-\alpha)^2 + \frac{a_2}{3}(z-\alpha)^3
+\ldots
\]
Durch Vergleich mit
\eqref{buch:ratint:rothsteintrager:eqn:integralform}
folgt, dass die gesuchten Koeffizienten $A_k$ die Residuen
des Integranden an den Nullstellen von $q(z)$ sein müssen.
\index{Residuum}%
Diese Koeffizienten lassen sich aber direkt aus den Polynomen 
$p(z)$ und $q(z)$ berechnen.

\begin{satz}
\label{buch:ratint:rothsteintrager:satz:residuum}
Sind $p(z)$ und $q(z)$ teilerfremd und $q(z)$ quadratfrei und
$\alpha$ eine Nullstelle des Nenners.
Dann ist das Residuum von $p(z)/q(z)$ an der Stelle $\alpha$
\[
A
=
\operatorname{res}\biggl(\alpha;\frac{p(z)}{q(z)}\biggr)
=
\frac{p(\alpha)}{q'(\alpha)}.
\]
\index{resz0f@$\operatorname{res}(z_0;f)$}%
\end{satz}

Der Satz ist eine direkte Folge des
Satzes~\ref{buch:analytisch:laurent:residuum:satz}.

